%!TEX root = ../main_es.tex
\appendix

\section{Apéndice técnico: BIDS}

El dataset se organiza en formato \textbf{Brain Imaging Data Structure (BIDS)}. Las señales originales se exportan desde BrainVision Analyzer y el procesamiento posterior se realiza en \Julia{}. La adaptación a BIDS busca trazabilidad, interoperabilidad y validación automática mediante BIDS Validator.

\begin{decisionmetodologica}
BIDS no se usa solo como convención de nombres. En este proyecto funciona como contrato de entrada: cada run debe poder resolverse desde sus metadatos, sus ficheros BrainVision enlazados y sus tablas de participantes/electrodos antes de entrar al pipeline.
\end{decisionmetodologica}

\subsection{Convenciones de carpetas y nombres}

\begin{tcolorbox}[
  enhanced,
  breakable,
  colback=NSGray,
  colframe=NSLine,
  boxrule=0.5pt,
  arc=1.5mm,
  title=\textbf{Reglas prácticas de nomenclatura BIDS}
]
\begin{itemize}[leftmargin=*,itemsep=2pt,topsep=2pt]
\item Los identificadores usan ceros iniciales: \bids{M5} $\to$ \bids{sub-M05}.
\item La jerarquía sigue \bids{sub-<id>/ses-<label>/eeg/}.
\item Las entidades BIDS (\texttt{sub-}, \texttt{ses-}, \texttt{task-}, \texttt{run-}) son minúsculas y separadas por guiones.
\item Los nombres usan ASCII, sin espacios ni acentos.
\item La condición ojos cerrados se codifica como \texttt{task-eyesclosed}; ojos abiertos como \texttt{task-eyesopen}.
\end{itemize}
\end{tcolorbox}

\begin{fullwidth}
\begin{figure}[H]
\centering
\resizebox{0.84\fulltextwidth}{!}{\NSBIDSBrainVision}
\caption{Resumen BIDS/BrainVision: tríptico de señal, metadatos tabulares y validación. Este esquema se repite aquí como referencia rápida del apéndice.}
\label{fig:appendix-bids-brainvision}
\end{figure}
\end{fullwidth}

\subsection{Tipos de archivo}

\begin{fullwidth}
\begin{table}[H]
\centering
\small
\caption{Tipos de archivo BIDS relevantes para el pipeline EEG.}
\label{tab:bids-files}
\begin{tabular}{@{}p{0.30\fulltextwidth}p{0.62\fulltextwidth}@{}}
\toprule
\textbf{Archivo o carpeta} & \textbf{Contenido y función} \\
\midrule
\bids{dataset\_description.json} & Metadatos generales del dataset: nombre, autores, versión BIDS y licencia. \\
\bids{participants.tsv} & Tabla de participantes: \texttt{participant\_id}, edad, sexo y grupo. \\
\bids{participants.json} & Descripción y unidades de las columnas de \bids{participants.tsv}. \\
\bids{README.md} & Documentación humana del dataset. \\
\bids{*\_eeg.eeg} & Señal EEG binaria continua en formato BrainVision. \\
\bids{*\_eeg.vhdr} & Cabecera BrainVision: ficheros asociados, canales y frecuencia de muestreo. Tras renombrar, \texttt{DataFile=} y \texttt{MarkerFile=} deben apuntar a los nuevos basenames. \\
\bids{*\_eeg.vmrk} & Marcadores BrainVision. Su campo \texttt{DataFile=} debe apuntar al \bids{.eeg} renombrado. \\
\bids{*\_eeg.json} & Sidecar BIDS-EEG: \texttt{SamplingFrequency}, \texttt{PowerLineFrequency}, referencia, tierra, número de canales y tipo de registro. \\
\bids{*\_events.tsv/json} & Eventos del run y descripción de sus columnas. \\
\bids{*\_electrodes.tsv} & Coordenadas y tipo de electrodos usados para visualización, QC y CSD. \\
\bottomrule
\end{tabular}
\end{table}
\end{fullwidth}

\subsection{Actualización de enlaces internos}

Después de renombrar los ficheros a basenames BIDS, las referencias internas de BrainVision deben actualizarse. En \bids{.vhdr}, \texttt{DataFile=} debe apuntar al nuevo \bids{.eeg} y \texttt{MarkerFile=} al nuevo \bids{.vmrk}; en \bids{.vmrk}, \texttt{DataFile=} debe apuntar al \bids{.eeg}. Sin esta actualización, herramientas compatibles con BrainVision y validadores BIDS pueden resolver mal los datos.

\section{Apéndice técnico: \Julia{}}

\Julia{} se utiliza como backend científico del proyecto porque combina sintaxis de alto nivel, buen rendimiento numérico y un ecosistema adecuado para matrices, señales, serialización y visualización. Frente a alternativas como Python o MATLAB, su principal ventaja práctica aquí es poder mantener código expresivo y cercano a la formulación matemática sin renunciar a rendimiento en bucles y operaciones de matriz \citep{bezanson2017julia}.

\subsection{Estructura del repositorio de análisis}

\begin{fullwidth}
\begin{table}[H]
\centering
\small
\caption{Mapa compacto del repositorio \Julia{} relacionado con el pipeline EEG.}
\label{tab:julia-repo-map}
\begin{tabular}{@{}p{0.18\fulltextwidth}p{0.72\fulltextwidth}@{}}
\toprule
\textbf{Directorio} & \textbf{Propósito} \\
\midrule
\texttt{config/} & Configuración de rutas y parámetros sin modificar código fuente. \\
\texttt{src/} & Módulos ejecutables: BIDS, IO, filtrado, ICA, procesamiento, FFT, CSD, wPLI, surrogates y FDR. \\
\texttt{script/} & Entradas ejecutables que orquestan módulos de \texttt{src/}. \\
\texttt{Pluto/} & Notebooks narrativos usados como documentación interactiva. \\
\texttt{data/} & Datos locales e intermedios por fase; no deberían versionarse como código. \\
\texttt{results/} & Figuras, tablas y logs regenerables por fase. \\
\texttt{docs/} & Exportaciones estáticas para documentación tipo GitHub Pages. \\
\bottomrule
\end{tabular}
\end{table}
\end{fullwidth}

\subsection{Figuras de resultados}

Las figuras de resultados del análisis (matrices de conectividad, espectros, comparaciones por grupo o condición) deben generarse desde \Julia{}, preferentemente en formatos vectoriales PDF/SVG cuando sea posible. Makie es adecuado para paneles complejos, matrices grandes y control fino de layouts; Plots.jl puede reservarse para exploración rápida.

\begin{lstlisting}[caption={Ejemplo mínimo de exportación de una matriz tipo conectividad en Julia.}]
using Plots
gr()
M = rand(8, 8)
p = plot(M, seriestype = :heatmap, xlabel = "Canal", ylabel = "Canal")
savefig(p, "connectivity_schematic.pdf")
\end{lstlisting}

\section{Apéndice técnico: dashboard NeuroMIND}

El dashboard de NeuroMIND es una interfaz web local para inspeccionar el estado del pipeline, revisar resultados por fase y conectar los artefactos generados por \Julia{} con decisiones de control de calidad. No sustituye al análisis estadístico ni al informe: funciona como capa de navegación, auditoría y depuración sobre \path{results/}.

\begin{controlqc}
\small
Este apéndice conserva únicamente la información necesaria para entender cómo el dashboard sostiene la trazabilidad del informe. Los detalles más operativos de mantenimiento, decisiones de carpeta, cambios recientes y pendientes de ejecución se registran en \path{CODEX.md} y \path{CHANGELOG.md}, no en el cuerpo científico principal.
\end{controlqc}

\begin{decisionmetodologica}
El dashboard debe interpretarse como herramienta de trazabilidad operativa. Sus paneles ayudan a detectar problemas de datos, configuración o ejecución, pero las conclusiones científicas deben formularse desde los outputs versionados, las tablas exportadas y el informe compilado.
\end{decisionmetodologica}

\subsection{Arquitectura}

La aplicación se lanza desde \path{scripts/launch_dashboard.jl}. El servidor local se implementa con Genie.jl en \path{src/webapp/App.jl} y sirve una interfaz HTML/JavaScript desde \path{web/views/dashboard.html}. Los endpoints leen configuración, datos BIDS y salidas ya generadas en \path{results/subjects/sub-<ID>/ses-<SES>/<task>/}.

\begin{fullwidth}
\begin{figure}[H]
\centering
\begin{tikzpicture}[
  font=\sffamily\small,
  node distance=1.0cm,
  box/.style={rounded corners=1.5mm, draw=NSLine, fill=NSGray, minimum width=3.0cm, minimum height=0.82cm, align=center},
  data/.style={box, fill=NSPetrol!8, draw=NSPetrol!70!black},
  code/.style={box, fill=NSBlue!7, draw=NSBlue!80!black},
  ui/.style={box, fill=NSOrange!9, draw=NSOrange!80!black},
  outputbox/.style={box, fill=NSPurple!8, draw=NSPurple!75!black},
  arr/.style={-{Stealth[length=2.2mm]}, thick, draw=NSPetrol!80!black}
]
  \node[data] (cfg) {config\\\path{*.toml}};
  \node[data, below=0.75cm of cfg] (bids) {data/BIDS\\raw + metadatos};
  \node[code, right=1.25cm of cfg] (pipeline) {pipeline Julia\\\path{src/} + \path{scripts/}};
  \node[outputbox, right=1.25cm of pipeline] (results) {outputs\\\texttt{results/subjects/...}};
  \node[code, below=0.75cm of results] (api) {Genie API\\\path{src/webapp/App.jl}};
  \node[ui, left=1.25cm of api] (html) {dashboard\\\path{web/views/dashboard.html}};
  \node[outputbox, below=0.75cm of api] (report) {informe\\\texttt{report/main\_es.tex}};

  \draw[arr] (cfg) -- (pipeline);
  \draw[arr] (bids) -- (pipeline);
  \draw[arr] (pipeline) -- (results);
  \draw[arr] (results) -- (api);
  \draw[arr] (api) -- (html);
  \draw[arr] (results) -- (report);
  \draw[arr, dashed, draw=NSOrange!90!black] (html) -- (report);
\end{tikzpicture}
\caption{Arquitectura operativa del dashboard: el pipeline genera outputs persistentes, la API Genie los expone al navegador y el informe reutiliza una selección estable de esos artefactos. La flecha discontinua indica que el dashboard guía la interpretación, pero no es la fuente primaria del texto científico.}
\label{fig:appendix-dashboard-architecture}
\end{figure}
\end{fullwidth}

\begin{fullwidth}
\begin{table}[H]
\centering
\small
\caption{Componentes principales del dashboard NeuroMIND.}
\label{tab:dashboard-components}
\begin{tabular}{@{}p{0.22\fulltextwidth}p{0.26\fulltextwidth}p{0.42\fulltextwidth}@{}}
\toprule
\textbf{Componente} & \textbf{Ruta} & \textbf{Función} \\
\midrule
Lanzador & \path{scripts/launch_dashboard.jl} & Activa el entorno \Julia{}, carga configuración y arranca el servidor local. \\
Servidor/API & \path{src/webapp/App.jl} & Define rutas HTTP, normaliza sujeto/sesión/condición y sirve CSV, JSON, imágenes y resúmenes. \\
Interfaz & \path{web/views/dashboard.html} & Renderiza paneles interactivos, tablas, gráficos canvas/SVG y acciones rápidas por fase. \\
Outputs & \path{results/subjects/...} & Fuente primaria de figuras, tablas, matrices, logs y snapshots usados por dashboard e informe. \\
Informe & \path{report/} & Selecciona resultados estables y los convierte en narrativa técnica reproducible. \\
\bottomrule
\end{tabular}
\end{table}
\end{fullwidth}

\subsection{Paneles y endpoints}

El dashboard usa una numeración más granular que la metodología del informe. La diferencia es intencionada: la metodología agrupa fases analíticas, mientras que el dashboard separa vistas de inspección para facilitar depuración rápida.

\begin{fullwidth}
\begin{table}[H]
\centering
\scriptsize
\caption{Mapa de paneles del dashboard, endpoints principales y uso previsto en el informe.}
\label{tab:dashboard-panels}
\begin{tabular}{@{}p{0.05\fulltextwidth}p{0.18\fulltextwidth}p{0.28\fulltextwidth}p{0.39\fulltextwidth}@{}}
\toprule
\textbf{Panel} & \textbf{Vista} & \textbf{Endpoint principal} & \textbf{Uso técnico} \\
\midrule
0 & Proyecto/dataset & \path{/api/project_overview} & Resumen de estudio, estado de datos y acciones iniciales. \\
1 & BIDS/metadatos & \path{/api/bids_metadata} & Verificación de estructura BIDS, canales, electrodos y metadatos. \\
2 & Señal raw & \path{/api/phase2_info} & Inspección de señal cruda, amplitud y espectro inicial. \\
3 & QC inicial & \path{/api/bids/channel_stats} & Estadísticos por canal, flags y topografía de calidad. \\
4 & Filtrado & \path{/api/filter_config}, \path{/api/filter_response} & Parámetros, cadena de filtros e impacto espectral. \\
5 & ICA & \path{/api/phase5_ica_info}, \path{/api/ica_features} & Componentes, scores de artefactos y revisión antes/después. \\
6 & Segmentación & \path{/api/phase6_segmentation} & Épocas, duración, cobertura y distribución temporal. \\
7 & Rechazo AR & \path{/api/phase7_ar} & Rechazo por amplitud/gradiente y ejemplos de épocas problemáticas. \\
8 & Espectral & \path{/api/phase8_spectral} & PSD, potencia por banda, regiones e índices espectrales. \\
9 & wPLI & \path{/api/phase9_wpli} & Heatmaps, red exploratoria, top aristas y resumen por banda. \\
10 & Surrogates/FDR & \path{/api/phase10_surrogates} & Observado, distribución nula, p/q-valores y conexiones significativas. \\
11 & Resultados finales & \path{/api/phase11_summary} & KPIs integrados y síntesis de ejecución. \\
12 & Exportación & \path{/api/phase12_files} & Inventario de archivos, categorías y rutas reproducibles. \\
13 & Transversal & \path{/api/phase13_transversal} & Comparación EM/control preparada para análisis grupal. \\
14 & Longitudinal & \path{/api/phase14_longitudinal} & Comparación T1/T2 preparada para seguimiento. \\
\bottomrule
\end{tabular}
\end{table}
\end{fullwidth}

\subsection{Relación con los resultados del informe}

Para el informe, el dashboard no debe exportarse como una colección de capturas. La estrategia preferible es usarlo como guía de selección: primero identifica artefactos útiles, después se incorporan al documento versiones estáticas, tablas resumidas o reconstrucciones vectoriales. En la ejecución \texttt{M05/T2/eyesclosed}, los paneles más relevantes para la sección de resultados son QC inicial, filtrado, ICA, segmentación/rechazo, análisis espectral, wPLI y surrogates/FDR.

\begin{controlqc}
Regla práctica: si un panel del dashboard responde a una decisión del pipeline, el informe debe citar el output persistente que sostiene esa decisión. Por ejemplo, un canal marcado como problemático debe trazarse a \path{qc_summary.csv}; una arista wPLI destacada debe trazarse a \path{connectivity_edges.csv} y, si se habla de inferencia, a \path{wpli_pvalues_*}, \path{wpli_qvalues_*} y \path{significant_connections.csv}.
\end{controlqc}

\subsection{Limitaciones}

\begin{itemize}[leftmargin=*,itemsep=2pt,topsep=2pt]
\item Los paneles interactivos facilitan inspección, pero no reemplazan la validación estadística ni la lectura metodológica.
\item Las vistas preparadas para análisis transversal y longitudinal requieren ejecuciones comparables por grupo, sesión y condición antes de interpretarse científicamente.
\item Las figuras del navegador pueden depender de escala de pantalla, selección de banda o umbrales interactivos; el informe debe usar artefactos exportados o figuras reconstruidas con parámetros explícitos.
\item La numeración del dashboard (0--14) no equivale uno a uno a la numeración metodológica del informe (0--8); por eso se documenta la correspondencia en la \Cref{tab:dashboard-panels}.
\end{itemize}

\section{Apéndice técnico: \LaTeX{}}

Este informe se compone con \LaTeX{} para separar contenido, estructura y presentación. La fuente textual permite versionar cambios, revisar diferencias y mantener referencias cruzadas, bibliografía, figuras y tablas de forma reproducible.

\begin{fullwidth}
\begin{table}[H]
\centering
\small
\caption{Componentes principales de la cadena \LaTeX{} del informe.}
\label{tab:latex-toolchain}
\begin{tabular}{@{}p{0.28\fulltextwidth}p{0.62\fulltextwidth}@{}}
\toprule
\textbf{Componente} & \textbf{Rol} \\
\midrule
\path{main_es.tex} & Orquestador de la versión española: paquetes, macros, portada e inclusión de secciones. \\
\path{sections_es/} & Contenido modular del cuaderno técnico. \\
\path{macros/neurosmart_style.tex} & Paleta, cajas destacadas, estilo de captions y títulos. \\
\path{assets/figures/} & Figuras organizadas por uso: \path{theory/}, \path{methodology/}, \path{connectivity/}, \path{results/}, \path{generated/} y \path{archive/}. \\
\texttt{refs.bib} & Base bibliográfica usada por \texttt{biber}. \\
\texttt{latexmk} & Compilación y reruns hasta resolver referencias y bibliografía. \\
\bottomrule
\end{tabular}
\end{table}
\end{fullwidth}

\subsection{Uso de \PGF{} \TikZ{}}

\TikZ{} se usa para diagramas conceptuales que conviene mantener como código: mapas de pipeline, BIDS, montaje, filtrado, ICA, segmentación, CSD/wPLI, surrogates/FDR y dashboard de outputs. Esto proporciona salida vectorial, tipografía coherente con el documento y edición trazable.

\begin{figure}[H]
\centering
\begin{tikzpicture}[
  scale=0.82,
  transform shape,
  node distance=1.2cm,
  channel/.style={circle, draw=NSBlue, fill=NSGray, minimum size=8pt, font=\scriptsize},
  edge/.style={-{Stealth[length=1.8pt]}, thin, NSPetrol}
]
  \node[channel] (Cz)  {Cz};
  \node[channel, above left=of Cz]  (Fz) {Fz};
  \node[channel, above right=of Cz] (Pz) {Pz};
  \node[channel, below=of Cz]       (Oz) {Oz};
  \draw[edge] (Fz) -- (Cz);
  \draw[edge] (Cz) -- (Pz);
  \draw[edge] (Cz) -- (Oz);
  \draw[edge] (Fz) to[bend left=18] (Pz);
\end{tikzpicture}
\caption{Esquema mínimo de conectividad sensorial en \TikZ{}.}
\label{fig:tikz-connectivity}
\end{figure}

\section{Apéndice técnico: \Git{} y \GitHub{}}

\Git{} se usa para registrar cambios del código \Julia{}, del informe \LaTeX{} y de los recursos pequeños necesarios para reproducir el trabajo. \GitHub{} añade copia remota, colaboración y publicación de documentación. Los datos raw, intermedios pesados y resultados regenerables deben excluirse o gestionarse explícitamente si alguna vez se decide versionarlos.

\begin{pendiente}
\small
La operación diaria de ramas, commits, incidencias y decisiones de mantenimiento se documenta fuera del cuerpo científico, principalmente en \path{CODEX.md}. Este apéndice solo fija el mínimo necesario para reproducibilidad y trazabilidad.
\end{pendiente}

\subsection{Flujo de trabajo mínimo}

\begin{fullwidth}
\begin{table}[H]
\centering
\small
\caption{Flujo básico con \Git{}/\GitHub{} para crear o clonar un repositorio.}
\label{tab:github-workflow}
\begin{tabular}{@{}p{0.22\fulltextwidth}p{0.30\fulltextwidth}p{0.38\fulltextwidth}@{}}
\toprule
\textbf{Escenario} & \textbf{Comando} & \textbf{Uso} \\
\midrule
Nuevo repositorio & \texttt{git init} & Inicializa un repositorio local. \\
 & \texttt{git remote add origin <url>} & Enlaza con GitHub. \\
 & \texttt{git add .} & Prepara cambios para commit. \\
 & \texttt{git commit -m "..."} & Crea una instantánea trazable. \\
 & \texttt{git push -u origin main} & Sube la rama principal. \\
\addlinespace
Repositorio existente & \texttt{git clone <url>} & Descarga una copia local. \\
 & \texttt{julia -e "using Pkg; Pkg.instantiate()"} & Reconstruye el entorno \Julia{} declarado. \\
 & \texttt{git pull} / \texttt{git push} & Sincroniza cambios. \\
\bottomrule
\end{tabular}
\end{table}
\end{fullwidth}

\subsection{Qué versionar}

\begin{tcolorbox}[
  enhanced,
  breakable,
  colback=NSGray,
  colframe=NSLine,
  boxrule=0.5pt,
  arc=1.5mm,
  title=\textbf{Contenido recomendado para commit}
]
\footnotesize
\begin{itemize}[leftmargin=*,itemsep=2pt,topsep=2pt]
\item \texttt{src/}, \texttt{script/} y \texttt{config/}.
\item \texttt{Project.toml} y, si se requiere reproducibilidad exacta, \texttt{Manifest.toml}.
\item \texttt{main\_es.tex}, \texttt{sections\_es/}, macros, bibliografía y figuras generadas pequeñas.
\item Documentación de estado como \texttt{README.md}, \texttt{CHANGELOG.md} y notas de ejecución.
\end{itemize}
\end{tcolorbox}

\section{Apéndice técnico: uso de herramientas de IA}

Este informe y el código asociado se han preparado con ayuda de herramientas de IA de forma transparente. Cursor se ha usado como entorno de edición asistida para navegar, escribir, refactorizar y depurar partes del proyecto. ChatGPT se ha usado para generar ideas, clarificar conceptos, revisar estructura y mejorar explicaciones. En ambos casos, las salidas se tratan como apoyo de trabajo: el contenido final debe revisarse, compilarse y contrastarse con el código y los datos disponibles.

Los detalles conversacionales, decisiones de edición y seguimiento de tareas no forman parte de la evidencia científica; se mantienen en documentos operativos del repositorio cuando aportan trazabilidad.
